% GENERATED FILE. Edit canonical lessons, then run npm run generate:books.
% canonical-source-hash: fnv1a64:9bda54dc438d76bd
% canonical-lessons: TE-C351-bandi, TE-C351-nirbandhincu, TE-C351-bandhincu, TE-C351-vidudala-ceyu, TE-C351-vidipincu

\chapter{Captives}
\label{ch:te-captives}
\hlchaptermodality{351}

\begin{chapteropening}
\textbf{By the end of this chapter:} \emph{I can say a captive, to detain, to lock up, to release and to free from captivity.}

Complete the last lesson of chapter 351: I can say a captive, to detain, to lock up, to release and to free from captivity.
\end{chapteropening}

\section[\texorpdfstring{\te{బందీ} (bandī)}{బందీ (bandī)}]{\te{బందీ} (bandī) — a captive; a hostage}

\label{lesson:TE-C351-bandi}



\begin{quote}
Before the new one: say the Telugu for an innocent person, then the Telugu for to punish.
\end{quote}

\subsection*{You'll want to know: \te{బందీ}}
\textbf{\te{బందీ}} — \emph{bandī} — \textquotedblleft{}a captive; a hostage\textquotedblright{}.

\begin{quote}
\te{దొంగలు} \te{ఇద్దరిని} \te{బందీలుగా} \te{పట్టుకున్నారు}.
\end{quote}

\begin{quote}
\emph{doṅgalu iddarini bandīlugā paṭṭukunnāru.}
\end{quote}

\begin{quote}
The robbers took two people captive.
\end{quote}

Plural: \te{బందీలు}, \emph{bandīlu}.

\te{బందీగా} \te{పట్టుకో}, to take captive.

\begin{cousinweb}
\textbf{\te{బందీ}} (\emph{bandī}) is a loan from Persian \emph{bandī}, \textquotedblleft{}a prisoner, a captive\textquotedblright{}, built on \emph{band}, \textquotedblleft{}a bond\textquotedblright{}.
\end{cousinweb}

\begin{grammarlens}[title={What you've built}]
One more word for telling what happened.
\end{grammarlens}

\subsection*{Your turn}
\begin{itemize}
\raggedright
  \item \emph{Say it:} \emph{bandī}
  \item \emph{Say it:} \emph{bandī}, once more
  \item \emph{Say it:} \emph{doṅgalu iddarini bandīlugā paṭṭukunnāru.} — The robbers took two people captive.
  \item \emph{Recall:} say the Telugu for an innocent person, then the Telugu for to punish, then say \emph{bandī} again
\end{itemize}

\begin{tcolorbox}[breakable,colback=teal!4,colframe=teal!35!black,title={Before you move on}]
What does \te{బందీ} mean? (\textquotedblleft{}A captive; a hostage\textquotedblright{}.) Say it once more.
\end{tcolorbox}
\section[\texorpdfstring{\te{నిర్బంధించు} (nirbandhiñcu)}{నిర్బంధించు (nirbandhiñcu)}]{\te{నిర్బంధించు} (nirbandhiñcu) — to detain; to confine}

\label{lesson:TE-C351-nirbandhincu}



\begin{quote}
Before the new one: say the Telugu for to punish, then the Telugu for a captive.
\end{quote}

\subsection*{You'll want to know: \te{నిర్బంధించు}}
\textbf{\te{నిర్బంధించు}} — \emph{nirbandhiñcu} — \textquotedblleft{}to detain; to confine\textquotedblright{}.

\begin{quote}
\te{పోలీసులు} \te{అనుమానితులను} \te{నిర్బంధించారు}.
\end{quote}

\begin{quote}
\emph{pōlīsulu anumānitulanu nirbandhiñcāru.}
\end{quote}

\begin{quote}
The police detained the suspects.
\end{quote}

Past (I …): \te{నిర్బంధించాను}, \emph{nirbandhiñcānu}.

An official word: police and governments detain.

\begin{cousinweb}
\textbf{\te{నిర్బంధించు}} (\emph{nirbandhiñcu}) is built on Sanskrit \emph{nirbandha}, \textquotedblleft{}perseverance, insistence\textquotedblright{} (CDIAL 7355), from \emph{bandh}, \textquotedblleft{}to bind\textquotedblright{}, with the Telugu verb ending \emph{-iñcu}.
\end{cousinweb}

\begin{grammarlens}[title={What you've built}]
One more word for telling what happened.
\end{grammarlens}

\subsection*{Your turn}
\begin{itemize}
\raggedright
  \item \emph{Say it:} \emph{nirbandhiñcu}
  \item \emph{Say it:} \emph{nirbandhiñcu}, once more
  \item \emph{Say it:} \emph{pōlīsulu anumānitulanu nirbandhiñcāru.} — The police detained the suspects.
  \item \emph{Recall:} say the Telugu for to punish, then the Telugu for a captive, then say \emph{nirbandhiñcu} again
\end{itemize}

\begin{tcolorbox}[breakable,colback=teal!4,colframe=teal!35!black,title={Before you move on}]
What does \te{నిర్బంధించు} mean? (\textquotedblleft{}To detain; to confine\textquotedblright{}.) Say it once more.
\end{tcolorbox}
\section[\texorpdfstring{\te{బంధించు} (bandhiñcu)}{బంధించు (bandhiñcu)}]{\te{బంధించు} (bandhiñcu) — to capture; to lock up}

\label{lesson:TE-C351-bandhincu}



\begin{quote}
Before the new one: say the Telugu for a captive, then the Telugu for to detain.
\end{quote}

\subsection*{You'll want to know: \te{బంధించు}}
\textbf{\te{బంధించు}} — \emph{bandhiñcu} — \textquotedblleft{}to capture; to lock up\textquotedblright{}.

\begin{quote}
\te{సైనికులు} \te{శత్రువులను} \te{ఒక} \te{గదిలో} \te{బంధించారు}.
\end{quote}

\begin{quote}
\emph{sainikulu śatruvulanu oka gadilō bandhiñcāru.}
\end{quote}

\begin{quote}
The soldiers locked the enemies up in a room.
\end{quote}

Past (I …): \te{బంధించాను}, \emph{bandhiñcānu}.

From \te{బంధం}, a bond: to bind someone.

\begin{cousinweb}
\textbf{\te{బంధించు}} (\emph{bandhiñcu}) is built on Sanskrit \emph{bandha}, \textquotedblleft{}a bond; custody\textquotedblright{} (CDIAL 9136), with the Telugu verb ending \emph{-iñcu}.
\end{cousinweb}

\begin{grammarlens}[title={What you've built}]
One more word for telling what happened.
\end{grammarlens}

\subsection*{Your turn}
\begin{itemize}
\raggedright
  \item \emph{Say it:} \emph{bandhiñcu}
  \item \emph{Say it:} \emph{bandhiñcu}, once more
  \item \emph{Say it:} \emph{sainikulu śatruvulanu oka gadilō bandhiñcāru.} — The soldiers locked the enemies up in a room.
  \item \emph{Recall:} say the Telugu for a captive, then the Telugu for to detain, then say \emph{bandhiñcu} again
\end{itemize}

\begin{tcolorbox}[breakable,colback=teal!4,colframe=teal!35!black,title={Before you move on}]
What does \te{బంధించు} mean? (\textquotedblleft{}To capture; to lock up\textquotedblright{}.) Say it once more.
\end{tcolorbox}
\section[\texorpdfstring{\te{విడుదల} \te{చేయు} (viḍudala cēyu)}{విడుదల చేయు (viḍudala cēyu)}]{\te{విడుదల} \te{చేయు} (viḍudala cēyu) — to release; to set free}

\label{lesson:TE-C351-vidudala-ceyu}



\begin{quote}
Before the new one: say the Telugu for to detain, then the Telugu for to capture.
\end{quote}

\subsection*{You'll want to know: \te{విడుదల} \te{చేయు}}
\textbf{\te{విడుదల} \te{చేయు}} — \emph{viḍudala cēyu} — \textquotedblleft{}to release; to set free\textquotedblright{}.

\begin{quote}
\te{పోలీసులు} \te{బందీలను} \te{విడుదల} \te{చేశారు}.
\end{quote}

\begin{quote}
\emph{pōlīsulu bandīlanu viḍudala cēśāru.}
\end{quote}

\begin{quote}
The police released the captives.
\end{quote}

Past (I …): \te{విడుదల} \te{చేశాను}, \emph{viḍudala cēśānu}.

\te{విడుదల} is also the release of a film or a book: \te{సినిమా} \te{విడుదల}.

\begin{cousinweb}
\textbf{\te{విడుదల} \te{చేయు}} (\emph{viḍudala cēyu}) is a native phrase: \emph{viḍudala}, \textquotedblleft{}release\textquotedblright{}, built on \emph{viḍu}, \textquotedblleft{}to loosen, let go\textquotedblright{} (DEDR 5393), with \emph{cēyu}, \textquotedblleft{}to do\textquotedblright{}.
\end{cousinweb}

\begin{grammarlens}[title={What you've built}]
One more word for telling what happened.
\end{grammarlens}

\subsection*{Your turn}
\begin{itemize}
\raggedright
  \item \emph{Say it:} \emph{viḍudala cēyu}
  \item \emph{Say it:} \emph{viḍudala cēyu}, once more
  \item \emph{Say it:} \emph{pōlīsulu bandīlanu viḍudala cēśāru.} — The police released the captives.
  \item \emph{Recall:} say the Telugu for to detain, then the Telugu for to capture, then say \emph{viḍudala cēyu} again
\end{itemize}

\begin{tcolorbox}[breakable,colback=teal!4,colframe=teal!35!black,title={Before you move on}]
What does \te{విడుదల} \te{చేయు} mean? (\textquotedblleft{}To release; to set free\textquotedblright{}.) Say it once more.
\end{tcolorbox}
\section[\texorpdfstring{\te{విడిపించు} (viḍipiñcu)}{విడిపించు (viḍipiñcu)}]{\te{విడిపించు} (viḍipiñcu) — to free from captivity}

\label{lesson:TE-C351-vidipincu}



\begin{quote}
Before the new one: say the Telugu for to capture, then the Telugu for to release.
\end{quote}

\subsection*{You'll want to know: \te{విడిపించు}}
\textbf{\te{విడిపించు}} — \emph{viḍipiñcu} — \textquotedblleft{}to free from captivity\textquotedblright{}.

\begin{quote}
\te{సైనికులు} \te{బందీలను} \te{క్షేమంగా} \te{విడిపించారు}.
\end{quote}

\begin{quote}
\emph{sainikulu bandīlanu kṣēmaṅgā viḍipiñcāru.}
\end{quote}

\begin{quote}
The soldiers freed the captives safely.
\end{quote}

Past (I …): \te{విడిపించాను}, \emph{viḍipiñcānu}.

The causative of \te{విడు}: to make someone come loose.

\begin{cousinweb}
\textbf{\te{విడిపించు}} (\emph{viḍipiñcu}) is native Dravidian: the causative of \emph{viḍu}, \textquotedblleft{}to be loosened, part with\textquotedblright{} (DEDR 5393).
\end{cousinweb}

\begin{grammarlens}[title={What you've built}]
One more word for telling what happened.
\end{grammarlens}

\subsection*{Your turn}
\begin{itemize}
\raggedright
  \item \emph{Say it:} \emph{viḍipiñcu}
  \item \emph{Say it:} \emph{viḍipiñcu}, once more
  \item \emph{Say it:} \emph{sainikulu bandīlanu kṣēmaṅgā viḍipiñcāru.} — The soldiers freed the captives safely.
  \item \emph{Recall:} say the Telugu for to capture, then the Telugu for to release, then say \emph{viḍipiñcu} again
\end{itemize}

\begin{tcolorbox}[breakable,colback=teal!4,colframe=teal!35!black,title={Before you move on}]
What does \te{విడిపించు} mean? (\textquotedblleft{}To free from captivity\textquotedblright{}.) Say it once more.
\end{tcolorbox}
